\scanpage{378}
\mspage{490}

\[
(29)\quad
\left\{\begin{aligned}
\underset{\substack{\Vert\,\mathrm{dfn.}\\\underline{u}\rho\underline{u}^{-1}}}{\underline{u}(\rho)}
 &=\underset{\substack{\Vert\\\operatorname{int}(\alpha)\rho}}{\underbrace{\xi\rho}},
 &\underline{u}(\sigma)&=\underset{\substack{\Vert\\\varepsilon\operatorname{int}(\beta)(\sigma)}}{\underbrace{\eta\sigma}},\\[1ex]
\det\underline{u}&=\mu
\end{aligned}\right.
\qquad\text{with }\xi=\alpha\rho(\alpha)^{-1},\quad
\eta=\varepsilon\beta\sigma(\beta)^{-1}
\]
\margin{watch the factor $\varepsilon$!}

\uline{Proposition} The morphism of schemes $\Theta$ is
\uline{characterized} by formulas (29), modulo a morphism
$\mathcal{G}_{\rho,\sigma}^{*}\xrightarrow{\lambda}\mu_2$.

Indeed, the first two formulas, which describe the action of
$\underline{u}$ on $\rho$ resp. $\sigma$, determine $\underline{u}$ modulo
$T_\rho$ resp. modulo $T_\sigma$, hence \inserted{(together)} modulo
$T_\rho\cap T_\sigma=\mathbb{G}_{\mathrm m}$, the third determines it
modulo $SG$, -- hence with the preceding ones, modulo $\mu_2$ -- i.e. another
morphism $\Theta'$ having the same properties
will be of the form
\[
\Theta'(g)=\lambda(g)\Theta(g)\qquad\text{where }\lambda:
\mathcal{G}_{\rho,\sigma}^{*}\longrightarrow\unsure{\mathcal{E}_2|\varepsilon_2}.
\]
\unsure{q.e.d.}

\uline{Remark} Let us note \inserted{moreover} that \deleted{\unsure{every}}
$\Theta$ is normalized by
\[
\Theta(\underbrace{1,1,1,1}_{1})=1
\]
which imposes \deleted{\unsure{an}}
\deleted{\unsure{additional}}\inserted{corresponding}
normalization $\lambda(1)=1$; on $\lambda$ -- but \unsure{unless I am mistaken} this
already implies (at least if $\rho,\sigma$ are $\frac12$ simple hence
$T_\rho,T_\sigma$ tori\ldots) that $\lambda$ is a \uline{homomorphism}
of group schemes, which here is \unsure{necessarily} trivial,
on \inserted{the component} $B'_0\times T_\rho\times T_\sigma$
(at least in characteristic $\ne2$) so we find essentially as the only
possible indeterminacy the multiplication \deleted{by} of
$\underline{u}$ by $\varepsilon$ \deleted{\unsure{all}}
(corresponding to the \inserted{canonical} homomorphism (nontrivial) from
$N'_\sigma/T_\sigma$ onto $\mu_2$)\ldots\ \deleted{\unsure{Error}}

\margin{If we worked \unsure{with} $S\mathcal{G}_{\rho,\sigma}$
instead of $\mathcal{G}_{\rho,\sigma}^{*}$, the indeterminacy would be
given by multiplication \unsure{by|$\mapsto$}
$\varepsilon^{i}\varepsilon'^{j}$
$(i,j\in\mathbb{Z}/2\mathbb{Z})$.}
\margin{377}

% Printed archive footer (not autograph).
\begin{center}
377\\
\texttt{https://grothendieck.umontpellier.fr}\\
Université de Montpellier, UMR CNRS 5149, Case courrier 051, Place Eugène Bataillon\\
34095 Montpellier cedex 5 - France
\end{center}

\scanpage{379}
\mspage{\unsure{490}}

(IX) \uline{The groups} $\mathcal{M}_{\rho,\sigma}$ \uline{and}
$S\mathcal{M}_{\rho,\sigma}$.

Let $\rho,\sigma$ be as always two sections of $G$, satisfying the
general conditions (cf. for example proposition page \unsure{478'|470'}), plus the
condition $\operatorname{Tr}\sigma=0$ (cf. \uline{V} -- for example lemma
page 455 and its corollary page 457) -- we could carry out their construction
also without this additional condition (contenting ourselves with
working with $\mathcal{G}_{\rho,\sigma}$ instead of
$\mathcal{G}_{\rho,\sigma}^{*}$) -- but \unsure{it matters little}
probably for what follows. We can if need be suppose
$\rho,\sigma$ normalized so that we have
$\operatorname{Tr}(\rho+\sigma)=1$.

We denote by $S\mathcal{M}_{\rho,\sigma}$ the product scheme
\[
(1)\qquad S\mathcal{M}_{\rho,\sigma}
=S\mathcal{G}_{\rho,\sigma}^{*}\times_S SG
\]
\margin{\unsure{In|For} what follows, the ``\unsure{officialization}''
\illegible{1} of the group \unsure{schemes|structures} $S_0\mathcal{M}_{\rho,\sigma}$,
cf. page 559 ff.}

Its sections are denoted
\[
(\alpha,\beta,\mu,\varepsilon,f)\ ;
\]
$f$ thus being an arbitrary section of $SG$
$(\simeq\mathrm{Sl}(2)_S)$.

In the case where $S=\operatorname{Spec}(A)$, $A=\widehat{\mathbb{Z}}$,
\[
G=\mathrm{Gl}(2)_S,\quad
\rho=\begin{pmatrix}0&1\\-1&1\end{pmatrix},\quad
\sigma=\begin{pmatrix}0&-1\\1&0\end{pmatrix},
\]
and where we thus have a canonical homomorphism
\[
(2)\qquad
\widehat{\unsure{\mathcal{G}|\mathcal{C}}}_{0,3}
\longrightarrow
G(\widehat{\mathbb{Z}})^{\prime\pm1}
\simeq\unsure{\mathrm{Gl}|\mathrm{Sl}}(2,\widehat{\mathbb{Z}})^{\pm1}/\{\pm1\}
\]
with values in the subgroup of \unsure{elements|elements} of determinant
$\pm1$,
\margin{378}

% Printed archive footer (not autograph).
\begin{center}
378\\
\texttt{https://grothendieck.umontpellier.fr}\\
Université de Montpellier, UMR CNRS 5149, Case courrier 051, Place Eugène Bataillon\\
34095 Montpellier cedex 5 - France
\end{center}

\scanpage{380}

sending $\widehat{\unsure{\mathcal{G}|\mathcal{C}}}_{0,3}^{+}$ into
$SG(\widehat{\mathbb{Z}})'=\mathrm{Sl}(2,\widehat{\mathbb{Z}})'
=\mathrm{Sl}(2,\widehat{\mathbb{Z}})/\unsure{\{\pm1\}|\pm1}$,
we are going to construct a variant of the group $\mathcal{M}_{0,3}^{\sim}$,
\inserted{(cf. page 390)}, which maps canonically into the product (1),
which will lead us to define a group structure on
$S\mathcal{M}_{\rho,\sigma}$, so as to make the said morphism
multiplicative.

Let us first consider the ``special'' variant
$S\mathcal{M}_{0,3}^{\sim}$ of $\mathcal{M}_{0,3}^{\sim}$, an extension
of $\mathcal{M}_{0,3}^{\sim}$ by $\widehat{\mathbb{Z}}^3$,
deduced from an extension of
\[
\unsure{\mathcal{A}|\mathcal{H}}_{0,3}^{\sim}
=\mathcal{M}_{0,3}^{\sim}/\widehat{\Pi}_{0,3}
\subset\unsure{\widehat{\widehat{\mathcal{T}}}_{0,3}|\widehat{\mathcal{T}}_{0,3}},
\]
\deleted{\unsure{p}} and even from an extension of
$\unsure{\widehat{\widehat{\mathcal{T}}}_{0,3}|\widehat{\mathcal{T}}_{0,3}}$
(via $\mathcal{M}_{0,3}^{\sim}\longrightarrow
\unsure{\widehat{\widehat{\mathcal{T}}}_{0,3}|\widehat{\mathcal{T}}_{0,3}}$),
consisting of the
\[
(3)\qquad
(u,\unsure{f_0|g_0},\unsure{f_1|g_1},\unsure{f_\infty|g_\infty})
\qquad
\left(u\in\unsure{\widehat{\widehat{\mathcal{T}}}_{0,3}|\widehat{\mathcal{T}}_{0,3}},\quad
g_0,g_1,g_\infty\in\unsure{\widehat{\Pi}_{0,3}}\right)
\]
satisfying
\[
(4)\qquad
u(\ell_i)=\operatorname{int}\bigl(\unsure{f|g}_{\dot{u}(i)}^{-1}\bigr)
\bigl(\ell_{\dot{u}(i)}^{\mu}\bigr)
\]
where $\mu$ is the multiplier of $u$, and $\dot{u}$ the action of $u$
on \inserted{the set} $\{0,1,\infty\}$ of loop classes.
The product is defined by
\[
\begin{aligned}
(5)\qquad &(u',f'_0,f'_1,f'_\infty)
(\deleted{\unsure{u}},f_0,f_1,f_\infty)\\
&=(u'u,\ f'_0u'(f_0),\ f'_1u'(f_1),\ \unsure{f'_\infty u'(f_\infty)})
\end{aligned}
\]
restricting ourselves (\unsure{deflation!}) to the case
$u,u'\in\mathcal{M}_{0,3}^{!\sim}$ i.e. $\dot{u},\dot{u}'=1$.

\margin{379}

% Printed archive footer (not autograph).
\begin{center}
379\\
\texttt{https://grothendieck.umontpellier.fr}\\
Université de Montpellier, UMR CNRS 5149, Case courrier 051, Place Eugène Bataillon\\
34095 Montpellier cedex 5 - France
\end{center}

\scanpage{381}
\mspage{491}

We have a surjective homomorphism $S\mathcal{M}_{0,3}^{\sim}\longrightarrow
\mathcal{M}_{0,3}^{\sim}$, whose kernel consists of the
\[
(6)\qquad (1,f_0,f_1,f_\infty)\qquad f_i\in L_i
\]
\unsure{thus we have}
\[
(7)\qquad
\diagram{\begin{tikzcd}[ampersand replacement=\&,column sep=large]
1\arrow[r]\&\widehat{\mathbb{Z}}^3\arrow[r]\&
S\mathcal{M}_{0,3}^{\sim}\arrow[r]\&
\mathcal{M}_{0,3}^{\sim}\arrow[r]\&1
\end{tikzcd}}
\]

We are going to define an interesting subgroup
\[
(8)\qquad S_0\mathcal{M}_{0,3}^{\sim}\subset
\mathcal{M}_{0,3}^{\sim},
\]
now using the definition of $\mathcal{M}_{0,3}^{\sim}$,
namely \inserted{\unsure{via}} the property of the
$u\in\mathcal{M}_{0,3}^{\sim}$ of commuting with
$\rho$ as outer actions, so that we will have
\[
(9)\qquad u\rho u^{-1}=g\rho\qquad\text{i.e. }[u,\rho]=g,
\quad g\in\widehat{\Pi}_{0,3}
\]

This allows us, when we have already \uline{chosen}
$f_i\in\widehat{\Pi}_{0,3}$ satisfying
\[
(10)\qquad u(\ell_i)=\operatorname{int}(f_i^{-1})(\ell_i^{\mu})
\]
(supposing for simplicity $u\in\mathcal{M}_{0,3}^{!\sim}$),
to deduce from it the \deleted{\unsure{forms|formulas}}
$f_{i+1},f_{i+2}$ via the formula
\[
\begin{aligned}
(11)\qquad f_{i+1}^{-1}&=g\rho(f_i^{-1})\,,\qquad
\inserted{\left[\begin{aligned}
\text{and }f_{i+2}^{-1}&=g\rho(f_{i+1}^{-1})\\
&=g\rho(g)\rho^2(f_i^{-1})
\end{aligned}\right]}\\[-1ex]
&\inserted{[\text{where }g=[u,\rho]\in\widehat{\Pi}_{0,3})}
\end{aligned}
\]
(we have set $\dot{\rho}(i)=i+1$ by abuse of language,
i.e. $0+1=1,\ 1+1=\infty,\ \infty+1=0$!) -- we
then \unsure{recover|will recover}
\[
f_i^{-1}=\unsure{f|g}_{(i+2)+1}^{-1}
=g\rho(f_{i+2}^{-1})
=\underbrace{g\rho(g)\rho^2(g)}_{1}\,
\underbrace{\rho^3(f_i^{-1})}_{f_i^{-1}}.
\]
\margin{380}

% Printed archive footer (not autograph).
\begin{center}
380\\
\texttt{https://grothendieck.umontpellier.fr}\\
Université de Montpellier, UMR CNRS 5149, Case courrier 051, Place Eugène Bataillon\\
34095 Montpellier cedex 5 - France
\end{center}

\scanpage{382}
\mspage{492}

The $(u,f_0,f_1,f_\infty)
\in S\mathcal{M}_{0,3}^{!\sim}$ which satisfy (11) for
$\forall\ i\in\{0,1,\infty\}$ (it suffices that it satisfy \unsure{it|them}
for two indices, for this to be satisfied for all
three) form a \uline{subgroup} of
$S\mathcal{M}_{0,3}^{!\sim}$, which I denote
$S_0\mathcal{M}_{0,3}^{!\sim}$ -- it also identifies with the group
of pairs $(u,f_0)$ $(u\in\mathcal{M},\ f_0\in\widehat{\Pi}_{0,3})$
satisfying
\[
(12)\qquad u(\ell_0)=\operatorname{int}(f_0^{-1})(\ell_0^{\mu})
\qquad(\inserted{\text{where }}\mu=\chi(u)).
\]

We have an extension structure
\[
(13)\qquad
\diagram{\begin{tikzcd}[ampersand replacement=\&,column sep=large]
1\arrow[r]\&\widehat{\mathbb{Z}}\arrow[r,no head]\&
S_0\mathcal{M}_{0,3}^{!\sim}\arrow[r]\&
\mathcal{M}_{0,3}^{!\sim}\arrow[r,no head]\&1
\end{tikzcd}},
\]
and a homomorphism of extensions
\[
(14)\qquad\left\{
\diagram{\begin{tikzcd}[ampersand replacement=\&,column sep=large,row sep=large]
1\arrow[r,no head]\&\widehat{\mathbb{Z}}\arrow[r,no head]\arrow[d,"\text{diagonal}"]\&
S_0\mathcal{M}_{0,3}^{!\sim}\arrow[r,no head]\arrow[d]\&
\mathcal{M}_{0,3}^{!\sim}\arrow[r,no head]\arrow[d]\arrow[d,no head,shift right=0.5ex]\&1\\
1\arrow[r,no head]\&\widehat{\mathbb{Z}}^3\arrow[r,no head]\&
S\mathcal{M}_{0,3}^{!\sim}\arrow[r,no head]\&
\mathcal{M}_{0,3}^{!\sim}\arrow[r,no head]\&1
\end{tikzcd}}\right.
\]
(where $S\mathcal{M}_{0,3}^{!\sim}$ is \inserted{defined as}
the \unsure{reverse|inverse} image of
$\mathcal{M}_{0,3}^{!\sim}\unsure{\subset|\in}\mathcal{M}_{0,3}^{\sim}$
in $S\mathcal{M}_{0,3}^{\sim}$). We can therefore say that
the extension $S\mathcal{M}_{0,3}^{!\sim}$ (of
$\mathcal{M}_{0,3}^{!\sim}$ by $\widehat{\mathbb{Z}}^3$) is
deduced, by the diagonal homomorphism
\[
(15)\qquad
\diagram{\begin{tikzcd}[ampersand replacement=\&,column sep=large,row sep=small]
\widehat{\mathbb{Z}}\arrow[r]\&\widehat{\mathbb{Z}}^3\\
\underline{n}\arrow[r,mapsto]\&
(\underline{n}_0=\underline{n},\ \underline{n}_1=\underline{n},\
\underline{n}_\infty=\underline{n})
\end{tikzcd}},
\]
from a ``finer'' extension $S_0\mathcal{M}_{0,3}^{!\sim}$ of
\margin{381}

% Printed archive footer (not autograph).
\begin{center}
381\\
\texttt{https://grothendieck.umontpellier.fr}\\
Université de Montpellier, UMR CNRS 5149, Case courrier 051, Place Eugène Bataillon\\
34095 Montpellier cedex 5 - France
\end{center}